\documentclass[11pt,a4paper]{article}

\usepackage[utf8]{inputenc}
\usepackage[T1]{fontenc}
\usepackage[brazil]{babel}
\usepackage[margin=2.5cm]{geometry}
\usepackage{amsmath,amssymb}
\usepackage{enumitem}
\usepackage{graphicx}
\usepackage{booktabs}
\usepackage{float}
\usepackage[hidelinks]{hyperref}

\graphicspath{{figures-panorama/}}
\setlength{\parskip}{0.4em}

\begin{document}

\begin{center}
  {\large\textbf{Panorama: comparando codificadores de fala}}\\[0.2em]
  {\small Explicabilidade acústica de detectores de \emph{deepfake} adaptados ao
  português brasileiro}
\end{center}

\section{O que foi feito nesta rodada}

\begin{itemize}[nosep]
  \item \textbf{Feito.} Rodamos o pipeline completo (desempenho, análise
        associativa, causal e confirmatória) de ponta a ponta.
  \item \textbf{Melhoramos.} Passamos de um único codificador para quatro,
        comparados sob a mesma receita; e produzimos figuras que os colocam lado a
        lado.
  \item \textbf{Adicionamos.} Uma lente nova de \emph{atenção temporal} (onde, no
        tempo, cada modelo olha), complementar às análises de frequência.
  \item \textbf{Aprofundamos.} Uma terceira lente de \emph{relevância} no espectro
        (DFT-LRP), que olha por dentro do modelo, agora conservativa nos três
        codificadores hf\_ssl (wav2vec2-base, hubert-base e wavlm-base).
  \item \textbf{Comparamos por dentro.} Colocamos a saliência dos três modelos no
        mesmo eixo, resumimos as faixas mais importantes de cada um numa tabela e
        medimos o quanto eles concordam sobre quais faixas importam.
  \item \textbf{Validamos.} Testamos a \emph{fidelidade} das faixas apontadas pelo
        DFT-LRP: removê-las e passar de novo pelo detector, comparando com um sorteio
        aleatório de faixas.
\end{itemize}

\section{A ideia, em poucas linhas}

Detectores de \emph{deepfake} de áudio de ponta costumam ser treinados em inglês.
Queremos entender quais pistas acústicas o detector usa para decidir se um áudio é
real (\emph{bonafide}) ou falso (\emph{spoof}), e como essas pistas mudam quando o
modelo é levemente adaptado ao português brasileiro.

Neste panorama damos um passo a mais: repetimos a mesma receita sobre quatro
codificadores de fala diferentes. Assim conseguimos ver se as respostas dependem
do codificador escolhido ou se há um padrão comum. A leitura é observacional, o
objetivo é descrever o comportamento sobre um conjunto representativo de áudios, de
forma verificável.

\section{O que comparamos}

Cada codificador entra congelado (sem treino no português). Sobre os
\emph{embeddings} que ele produz, treinamos um \emph{head} leve (regressão
logística) para o português; esse head é a ``adaptação''. Os quatro codificadores
comparados:

\begin{itemize}[nosep]
  \item \textbf{XLS-R} (multilíngue, cerca de 300 milhões de parâmetros, com
        \emph{fine-tuning} anti-\emph{deepfake});
  \item \textbf{wav2vec2-base} (cerca de 95 milhões);
  \item \textbf{hubert-base-ls960};
  \item \textbf{wavlm-base}.
\end{itemize}

Todos são avaliados no mesmo conjunto de análise, na mesma grade de 24 faixas de
frequência (escala mel) e com as mesmas análises: desempenho, uma análise
\emph{associativa} (quais faixas acompanham a decisão), uma análise \emph{causal}
(quais faixas, ao serem removidas, de fato mudam a decisão) e uma etapa
confirmatória ligada aos erros.

\section{As perguntas e as hipóteses}

\begin{enumerate}[label=\textbf{Q\arabic*.}]
  \item Em quais faixas de frequência cada detector adaptado se apoia, e há um padrão
        em comum entre os codificadores?
  \item Essas faixas de fato sustentam a decisão (efeito causal e fidelidade) e se
        relacionam com os erros?
\end{enumerate}

\begin{itemize}[nosep]
  \item \textbf{H1 (descritiva).} Cada codificador se apoia num conjunto próprio de
        faixas de frequência, e num sentido próprio (rumo a \emph{spoof} ou a
        \emph{bonafide}); observamos se há alguma estrutura em comum.
  \item \textbf{H2 (causal).} Algumas faixas, ao serem removidas do áudio, de fato
        mudam a decisão; essas faixas devem sustentar a decisão de forma verificável
        (fidelidade), não apenas acompanhá-la.
  \item \textbf{H3 (inferencial).} As pistas mais importantes se comportam de forma
        diferente quando o detector acerta e quando erra.
\end{itemize}

\section{Panorama dos achados}

\subsection*{Desempenho: todos adaptam, o XLS-R lidera}

A métrica principal é a EER (\emph{equal error rate}): o ponto em que a taxa de
deixar passar um \emph{deepfake} se iguala à de rejeitar um áudio real. É um resumo
do detector em um número só, e \emph{quanto menor, melhor}. A
Figura~\ref{fig:eer} e a Tabela~\ref{tab:perf} mostram os quatro adaptados.

\begin{table}[H]
  \centering
  \small
  \begin{tabular}{lccc}
    \toprule
    Codificador & EER adaptado (\%) & MCC & Acurácia \\
    \midrule
    XLS-R              & 20{,}9 & 0{,}58 & 0{,}79 \\
    wav2vec2-base      & 22{,}1 & 0{,}56 & 0{,}78 \\
    hubert-base-ls960  & 26{,}2 & 0{,}47 & 0{,}74 \\
    wavlm-base         & 29{,}0 & 0{,}42 & 0{,}71 \\
    \bottomrule
  \end{tabular}
  \caption{Desempenho do detector adaptado por codificador. Só o XLS-R traz também
  um head original (\emph{zero-shot}), com EER de 25{,}8\%, que a adaptação reduz
  para 20{,}9\%.}
  \label{tab:perf}
\end{table}

Todos os codificadores atingem uma EER entre cerca de 21\% e 29\%. O XLS-R, maior,
multilíngue e já ajustado para a tarefa, tem a melhor detecção; os três modelos
\emph{base}, menores, ficam atrás, com o wavlm por último. É um resultado esperado:
mais capacidade, pré-treino multilíngue e adaptação de fábrica ajudam. Esse
desempenho é o ponto de partida do estudo, pois só faz sentido investigar as pistas de
detectores que de fato detectam.

\begin{figure}[H]
  \centering
  \includegraphics[width=0.72\textwidth]{encoder_comparison_eer}
  \caption{EER do detector adaptado por codificador (quanto menor, melhor).}
  \label{fig:eer}
\end{figure}

\subsection*{H1: a ``impressão digital'' associativa muda muito entre codificadores}

Na Figura~\ref{fig:assoc}, cada célula mede a associação entre a energia de uma
faixa de frequência e a decisão do detector (correlação de Spearman, calculada
dentro de cada classe para evitar correlação espúria). A cor indica o sentido:
vermelho quando mais energia naquela faixa acompanha um score mais alto de
\emph{spoof} (\(\rho>0\)); azul quando mais energia acompanha um score mais baixo,
na direção de \emph{bonafide} (\(\rho<0\)). O tom indica a intensidade.

A leitura é que a ``impressão digital'' associativa depende bastante do
codificador. O XLS-R e o wav2vec2-base mostram um padrão estruturado, com inversão
de sentido ao longo do espectro (graves em vermelho, ou seja associados a
\emph{spoof}, e faixas médias em azul, associadas a \emph{bonafide}). Já o hubert e o
wavlm ficam espalhados, com associações fracas e quase todas no mesmo sentido (vermelho
suave). É o registro descritivo de H1: as faixas que acompanham a decisão, e o
sentido dessa associação, não são os mesmos entre os modelos.

\begin{figure}[H]
  \centering
  \includegraphics[width=\textwidth]{encoder_assoc_heatmap}
  \caption{Perfil associativo (Spearman com sinal) por faixa, um codificador por
  linha. Vermelho: mais energia acompanha \emph{spoof} (\(\rho>0\)); azul: mais
  energia acompanha \emph{bonafide} (\(\rho<0\)). Registro descritivo (H1).}
  \label{fig:assoc}
\end{figure}

\subsection*{H2: os graves concentram o efeito causal}

Aqui intervimos no áudio: removemos uma faixa de frequência por vez (filtro de fase
zero) e medimos quanto a probabilidade de \emph{spoof} muda
(Figura~\ref{fig:occ}). Valores positivos indicam faixas que \emph{sustentam} a
decisão de falso (removê-las derruba o score de \emph{spoof}); valores negativos
indicam faixas que funcionavam como indício de \emph{real} (removê-las aumenta o
score de \emph{spoof}). A região sombreada é o intervalo de confiança de 95\%.

O padrão comum é que as \emph{frequências graves} (abaixo de cerca de 300~Hz)
concentram o maior efeito, enquanto acima de 1~kHz os efeitos são pequenos e
parecidos entre os modelos. O sentido nos graves, porém, varia: para o wavlm e o
hubert as faixas mais graves funcionam como indício de \emph{real} (removê-las
aumenta o score de \emph{spoof}), e o wav2vec2 tem um mergulho marcante perto de
200--300~Hz. É a evidência causal de H2, a parte mais forte do estudo, por ser a
única baseada em intervenção direta.

\begin{figure}[H]
  \centering
  \includegraphics[width=0.72\textwidth]{encoder_occlusion_overlay}
  \caption{Perfil causal por faixa (uma linha por codificador). Positivo: a faixa
  sustenta a decisão de \emph{spoof}; negativo: a faixa era indício de
  \emph{bonafide}. Sombra: intervalo de confiança de 95\%. Registro causal (H2).}
  \label{fig:occ}
\end{figure}

\subsection*{As duas lentes apontam para as mesmas faixas?}

Vale checar se a faixa mais \emph{associada} à decisão (H1) é também a que mais a
\emph{muda} quando removida (H2). A Figura~\ref{fig:conv} resume isso por um único
número por codificador, com sinal: positivo quando associação e efeito causal
apontam juntos, próximo de zero ou negativo quando não.

Com cautela: os valores ficam próximos de zero ou negativos nos quatro modelos
(por exemplo, \(-0{,}20\) no XLS-R e \(-0{,}39\) no wavlm). Ou seja, as duas lentes
ainda \emph{não} apontam de forma limpa para as mesmas faixas. Uma causa provável é
um efeito colateral da própria oclusão: ao remover uma faixa, mexemos também na
energia total do clipe, o que pode confundir o efeito daquela faixa específica.
Por isso tratamos a análise causal como a evidência mais forte e a associativa como
descritiva, e deixamos o controle desse efeito colateral como próximo passo.

\begin{figure}[H]
  \centering
  \includegraphics[width=0.72\textwidth]{encoder_convergence_forest}
  \caption{Convergência entre a análise associativa (H1) e a causal (H2), um ponto
  por codificador com intervalo de confiança. Positivo indicaria que as faixas mais
  associadas são também as mais causais; os valores observados ficam perto de zero
  ou negativos.}
  \label{fig:conv}
\end{figure}

\subsection*{Uma lente extra: onde, no tempo, cada modelo olha}

Além do espectro, olhamos o eixo do \emph{tempo}. A Figura~\ref{fig:roll} usa o
\emph{attention roll-out}, que consolida a atenção interna do codificador ao longo
das camadas e mostra quais instantes do áudio ele mais atende. Usamos os mesmos
quatro clipes nos três codificadores (um por coluna), para uma comparação justa.

A leitura é qualitativa. Nos três modelos a atenção acompanha os trechos com voz e
cai nos silêncios, o que é um bom sinal de sanidade. A grade também deixa ver onde
os modelos \emph{discordam}: na terceira linha, um áudio real aceito pelo wavlm e
pelo hubert é marcado como falso pelo wav2vec2 (o P(\emph{spoof}) de cada célula é a
decisão daquele modelo).

\begin{figure}[H]
  \centering
  \includegraphics[width=\textwidth]{attention_rollout_compare}
  \caption{Atenção temporal (\emph{roll-out}) dos codificadores nos mesmos clipes.
  Linhas: clipes (quadrantes de acerto/erro do detector de referência). Colunas:
  codificadores atendendo ao \emph{mesmo} áudio. \textbf{Ressalva:} é a atenção
  temporal do codificador congelado, não a decisão do \emph{head} logístico e não é
  específica de frequência. Vista complementar e ilustrativa.}
  \label{fig:roll}
\end{figure}

\section{Uma lente por dentro do modelo (DFT-LRP)}

As duas lentes anteriores olham o áudio de fora (correlação em H1, remoção de
faixas em H2). Aqui abrimos o modelo. Usamos a \emph{relevância} (LRP, sigla de
\emph{Layer-wise Relevance Propagation}): partindo da decisão, distribuímos a
contribuição de volta por cada camada até a forma de onda, e uma ponte matemática
(DFT-LRP) leva essa contribuição do tempo para a frequência. A relevância é a
parcela do \emph{logit} de \emph{spoof} (a nota bruta antes de virar probabilidade)
que cada frequência carrega, positiva quando empurra para \emph{spoof} e negativa
quando empurra para \emph{bonafide}.

Esta lente depende das regras internas de atenção do codificador. Nós as
implementamos de forma \emph{conservativa} (a soma pelas frequências recupera a
decisão, checado por um teste numérico) para os três codificadores hf\_ssl:
wav2vec2-base e hubert-base, que compartilham a atenção padrão, e o
wavlm-base, que usa um tipo de atenção diferente e por isso recebeu uma regra
dedicada. Todos passam no teste (resíduo da ordem de \(10^{-4}\) ou menor). Uma
ressalva vale para o wavlm: nele os pesos internos carregam uma parcela muito grande
da nota, então a relevância atribuída à forma de onda fica com valores inflados. Por
isso, comparamos os perfis pela forma e pelo sentido, não pelo valor absoluto.

A Figura~\ref{fig:lrp-overlay} sobrepõe o perfil de relevância por faixa dos três
codificadores hf\_ssl. O XLS-R original usa outra biblioteca (fairseq) e fica fora
desta lente. A leitura é cautelosa: os intervalos de confiança são
largos e, pela ressalva acima, a magnitude do wavlm não é diretamente comparável às
demais. Ainda assim, o padrão ecoa a análise causal (H2): a relevância se concentra
abaixo de cerca de 1~kHz e fica próxima de zero nos agudos. O wavlm mostra um pico
positivo forte (rumo a \emph{spoof}) nos graves, perto de 130--500~Hz; o hubert tem
um pico positivo mais modesto perto de 130--150~Hz; o wav2vec2 tem uma queda (rumo a
\emph{bonafide}) perto de 1{,}3--1{,}5~kHz. É um indício
adicional, por dentro do modelo, de que a ação está nos graves e médios-graves.

\begin{figure}[H]
  \centering
  \includegraphics[width=0.82\textwidth]{encoder_lrp_overlay}
  \caption{Relevância por faixa (DFT-LRP) do detector adaptado, uma linha por
  codificador. Positivo: a faixa empurra a decisão para \emph{spoof}; negativo: para
  \emph{bonafide}. Os três usam o método conservativo (linha cheia). No wavlm, ler
  pela forma e pelo sentido, não pela magnitude (os vieses internos carregam grande
  parte do \emph{logit}). Sombra: intervalo de confiança de 95\%.}
  \label{fig:lrp-overlay}
\end{figure}

Descendo ao nível de clipe, a Figura~\ref{fig:lrp-sal} mostra \emph{onde}, no
espectro, o wav2vec2-base mais se apoia. É a relevância em módulo (sem distinguir o
sentido), somada e normalizada por clipe, então lê-se como a fração da atenção que
cai em cada frequência (em \%). O
foco se concentra em torno de 1{,}4~kHz (cerca de 13\% da atenção), com um segundo
pico nos graves perto de 150~Hz (região da frequência fundamental da voz). Como é
magnitude, parte desse pico pode vir do ganho da própria primeira etapa do modelo (o
\emph{front-end} convolucional) nessa banda, então lemos como concentração de atenção,
não como prova de causa.

\begin{figure}[H]
  \centering
  \includegraphics[width=0.72\textwidth]{dft_lrp_salience}
  \caption{Saliência em frequência (DFT-LRP) do detector adaptado sobre o
  wav2vec2-base. Altura: fração da relevância em módulo em cada frequência (\% por
  bin), média entre clipes; sombra: intervalo de confiança de 95\%.
  Sem distinguir o sentido, indica onde o modelo olha, não para qual lado. O foco se
  concentra em torno de 1{,}4~kHz.}
  \label{fig:lrp-sal}
\end{figure}

\subsection{Onde cada modelo olha, lado a lado}

A Figura~\ref{fig:sal-overlay} coloca essa mesma saliência dos três codificadores no
mesmo eixo. Fica claro que o foco em frequência não é o mesmo entre eles. O
wav2vec2-base tem um pico marcante e isolado perto de 1{,}4~kHz; o hubert-base e o
wavlm-base concentram a atenção nos graves (pico perto de 220--250~Hz) e caem de
forma suave depois disso. É a mesma leitura da lente causal e das médias populacionais
em frequência: cada modelo se apoia numa sub-região diferente do espectro, com os
graves em comum para hubert e wavlm.

\begin{figure}[H]
  \centering
  \includegraphics[width=0.9\textwidth]{encoder_lrp_salience_overlay}
  \caption{Saliência em frequência (DFT-LRP) sobreposta, uma linha por codificador.
  Altura: fração da relevância em módulo por frequência (\% por bin), média entre
  clipes. Sem distinguir o sentido: indica onde cada modelo se apoia, não para qual lado.
  Rótulo: frequência de pico por codificador.}
  \label{fig:sal-overlay}
\end{figure}

A Tabela (Figura~\ref{fig:top-bands}) resume o mesmo de forma direta: as cinco faixas
mais importantes de cada modelo, ordenadas pela relevância em módulo. A faixa
\emph{Top~1} confirma o contraste: 1{,}43~kHz no wav2vec2, contra 59~Hz no hubert e
142~Hz no wavlm. Nenhuma faixa é \emph{Top~1} em mais de um modelo.

\begin{figure}[H]
  \centering
  \includegraphics[width=0.95\textwidth]{encoder_lrp_top_bands}
  \caption{Faixas mais importantes por codificador (DFT-LRP), ordenadas pela
  relevância em módulo. A linha \emph{Top~1} está destacada. Leitura descritiva: cada
  modelo prioriza faixas diferentes.}
  \label{fig:top-bands}
\end{figure}

A mesma relevância pode ser vista no plano tempo-frequência de clipes isolados
(STDFT-LRP). A Figura~\ref{fig:stdft} traz um exemplo predito \emph{spoof} e um
predito \emph{bonafide}. O mapa central colore cada instante e frequência pela
relevância (vermelho para \emph{spoof}, azul para \emph{bonafide}); o painel à
esquerda soma no tempo (relevância por frequência) e o de baixo mostra a forma de
onda. A relevância acompanha os trechos com voz e a estrutura das harmônicas; no
exemplo \emph{spoof} há um risco vertical marcante (um som breve e abrupto) puxando
para \emph{spoof}.

\begin{figure}[H]
  \centering
  \includegraphics[width=0.49\textwidth]{stdft_spoof}
  \includegraphics[width=0.49\textwidth]{stdft_bonafide}
  \caption{STDFT-LRP em dois clipes (wav2vec2-base). Esquerda: predito \emph{spoof};
  direita: predito \emph{bonafide}. Mapa central: relevância por tempo e frequência
  (vermelho empurra para \emph{spoof}, azul para \emph{bonafide}). Painel lateral:
  relevância somada no tempo, por frequência. Painel inferior: forma de onda.
  Ilustrativo, um clipe por lado.}
  \label{fig:stdft}
\end{figure}

\subsection{Visão populacional (média sobre muitos clipes)}

Os mapas acima são de clipes isolados. Para uma visão de conjunto, calculamos a
relevância STDFT-LRP de cada áudio do conjunto de teste (2820 clipes: 1320
\emph{bonafide} e 1500 \emph{spoof}) e tiramos a média. Isso é possível porque todos
os clipes têm o mesmo comprimento fixo, então a grade do STFT é idêntica entre eles.

Aqui há um cuidado importante. O eixo do tempo não está alinhado entre enunciados
(cada áudio fala em instantes diferentes, com silêncios e durações distintos), então
uma média instante a instante dependeria de como os clipes se arranjam no conjunto, e
não de uma propriedade do modelo. Por isso \emph{eliminamos o tempo}: para cada
frequência, tiramos a média da relevância ao longo de todo o clipe. O que resta é
comparável entre modelos e não depende do arranjo temporal do dataset.

A Figura~\ref{fig:pop-mag} mostra essa magnitude da relevância (sem distinguir o
sentido) já colapsada no tempo, num recorte de 0 a 1~kHz, faixa onde a atenção se
concentra. Cada coluna é um codificador, normalizada pelo seu próprio máximo. Os três
se apoiam nas frequências graves, cada um numa sub-faixa um pouco diferente: o
wav2vec2 perto de 150~Hz e o hubert e o wavlm um pouco acima (perto de 220--270~Hz).

\begin{figure}[H]
  \centering
  \includegraphics[width=0.7\textwidth]{stdft_lrp_population_strip_magnitude_1khz}
  \caption{Importância em frequência populacional (STDFT-LRP), com o tempo eliminado
  por média, recorte de 0 a 1~kHz. Cada coluna é um codificador; a cor é a magnitude
  média da relevância por frequência, normalizada pelo máximo de cada modelo. Mais
  claro: frequência em que o modelo mais se apoia. Sem eixo de tempo, então a leitura
  não depende do arranjo dos clipes no dataset.}
  \label{fig:pop-mag}
\end{figure}

A Figura~\ref{fig:pop-freq} resume a mesma informação em uma dimensão: a média no
tempo, por frequência, normalizada para somar 100\% no espectro (uma fração de
atenção, como o mapa de saliência). Os picos são coerentes com a leitura anterior:
graves para os três e um segundo
pico perto de 1{,}4~kHz apenas no wav2vec2. Essa
curva usa um caminho de cálculo diferente da Figura~\ref{fig:lrp-sal} (STFT curto,
não a transformada do clipe inteiro) e chega ao mesmo padrão, o que serve de
verificação cruzada.

\begin{figure}[H]
  \centering
  \includegraphics[width=0.82\textwidth]{stdft_lrp_population_freq_magnitude}
  \caption{Curva em frequência da atenção populacional (STDFT-LRP): média no tempo
  da magnitude da relevância, por frequência, uma linha por codificador. Normalizada
  para somar 100\% no espectro (fração de atenção por bin). Sem distinguir o sentido: indica
  onde cada modelo se apoia, não para qual lado. Rótulos: frequência de pico por
  codificador.}
  \label{fig:pop-freq}
\end{figure}

\subsection{Quanto os modelos concordam sobre quais faixas importam?}

As figuras acima sugerem que os modelos priorizam faixas diferentes. Para colocar um
número nisso, a Figura~\ref{fig:lrp-corr} compara os perfis de relevância por faixa
dois a dois com a correlação de Spearman (correlação entre as ordens das 24 faixas).
O sentido da leitura: \(\rho\) próximo de \(+1\) indica que os dois modelos ordenam
as faixas de forma parecida (concordam sobre quais importam mais); próximo de zero
indica pouca relação; negativo indicaria ordens invertidas.

A concordância observada é, no máximo, fraca a moderada e positiva: hubert e wavlm
são os mais parecidos (\(\rho = 0{,}35\)), wav2vec2 e hubert vêm depois
(\(\rho = 0{,}24\)), e wav2vec2 e wavlm são praticamente independentes
(\(\rho = -0{,}03\)). Ou seja, não há uma faixa universal que todos usem da mesma
forma. Uma ressalva importante: cada \(\rho\) é calculado sobre apenas 24 faixas, o
que o torna sugestivo, não conclusivo. A unidade natural de incerteza é o número de
clipes, então uma versão futura com intervalo de confiança por reamostragem dos
clipes deixaria essa leitura mais firme.

\begin{figure}[H]
  \centering
  \includegraphics[width=0.7\textwidth]{encoder_lrp_correlation}
  \caption{Concordância entre codificadores: correlação de Spearman entre os perfis
  de relevância por faixa (DFT-LRP), dois a dois. Vermelho forte: ordenam as faixas de
  forma parecida (\(\rho>0\)); tons claros: pouca relação. Registro descritivo sobre
  as 24 faixas (sugestivo, ainda sem intervalo de confiança).}
  \label{fig:lrp-corr}
\end{figure}

\subsection{As faixas apontadas de fato sustentam a decisão? (fidelidade)}

As lentes anteriores dizem \emph{onde} o modelo se apoia. Falta uma pergunta: essas
faixas são mesmo as que decidem, ou só acompanham a decisão por acaso? Para checar,
usamos o teste de \textbf{necessidade} (comprehensiveness): removemos do áudio as faixas
apontadas pelo DFT-LRP e passamos de novo pelo detector. Se a decisão cai, essas faixas
são \emph{necessárias}.

Comparamos com um \emph{sorteio aleatório}: remover o mesmo número de faixas escolhidas
ao acaso. A ideia é direta: se as faixas do DFT-LRP importam mais que faixas quaisquer,
removê-las deve derrubar a decisão mais do que o sorteio. Repetimos isso para vários
números de faixas e resumimos num valor médio (chamado AOPC, a área sob essa curva); a
barra de erro é a incerteza (intervalo de confiança de 95\%, obtido reamostrando os
clipes).

Como ler a Figura~\ref{fig:faith}: são dois painéis, um por classe (\emph{spoof} e
\emph{bonafide}). Em cada painel aparecem os três codificadores, e para cada um duas
barras: a cheia (remover as faixas do DFT-LRP) e a hachurada (remover faixas ao acaso).
Quanto mais alta a barra, mais a decisão cai. A barra cheia acima da hachurada é o sinal
de que a explicação é fiel.

O resultado é firme: remover as faixas do DFT-LRP derruba a decisão bem mais do que
remover faixas ao acaso, nos três modelos e nas duas classes (só o wavlm em \emph{spoof}
fica no limite). Isso indica que as faixas apontadas pesam de verdade, não são só
correlação, o que reforça a leitura de H2.

\begin{figure}[H]
  \centering
  \includegraphics[width=\textwidth]{faithfulness_compare_comprehensiveness}
  \caption[Necessidade das faixas (comprehensiveness)]{Necessidade das faixas
  (comprehensiveness): remover as faixas do DFT-LRP contra remover faixas ao acaso, um
  grupo de barras por codificador, um painel por classe. Amostra de 500 clipes por
  classe.\\[3pt]
  \begin{minipage}{0.95\linewidth}
  \begin{itemize}[nosep,leftmargin=1.5em,topsep=0pt]
    \item Altura da barra: AOPC (o quanto a decisão cai, em média, ao remover as faixas);
          barra de erro: intervalo de confiança de 95\%.
    \item Mais alto = a decisão cai mais, ou seja, as faixas removidas eram mais
          necessárias.
    \item Barra cheia (DFT-LRP) acima da hachurada (aleatória) = sinal de fidelidade.
  \end{itemize}
  \end{minipage}}
  \label{fig:faith}
\end{figure}

\section{Panorama em uma frase}

Dá para adaptar detectores de \emph{deepfake} treinados em inglês ao português com
um head leve (EER de 21\% a 29\%, XLS-R à frente). As pistas que cada codificador
associa e usa causalmente diferem bastante, com um ponto em comum: as frequências
graves concentram o efeito causal. As análises associativa e causal ainda não
convergem de forma limpa, o que fica como próximo passo (controlar o efeito
colateral da oclusão). A atenção temporal entra como uma lente complementar, útil
para inspecionar casos e discordâncias. Por dentro do modelo, uma terceira lente
(DFT-LRP), agora conservativa nos três codificadores hf\_ssl, volta a apontar os
graves e médios-graves, coerente com a análise causal; comparados lado a lado, os
três modelos concordam apenas de forma fraca a moderada sobre quais faixas importam
(Spearman de \(-0{,}03\) a \(0{,}35\)), sem uma faixa universal. Um teste de fidelidade
confirma que as faixas apontadas pelo DFT-LRP são necessárias à decisão nos três
modelos: removê-las derruba a decisão mais que remover faixas aleatórias.

\end{document}
